%! The Normal Distribution and z-Scores — Unit 1, Lesson 1.7 — Guided Notes & Practice (Answer Key)
% MIRROR of ../notes/main.tex, differing in exactly FIVE ways: apstats-key in place of
% apstats-boxes, the \pageheader's second argument, \termblank -> \termans, \blank -> \ans, and
% the fourth argument of every \pcell. Every `work` block, TikZ block, height, width, sub-row
% seam and \boxguard is byte-identical, and every prose cut is made in BOTH files — keep it
% that way.
% THE in-class document, in the MAIN IDEAS / NOTES shape (see LESSON_SHAPE.md §1): the
% vocabulary box, then ONE two-column guidednotes table. Each row is one idea — a label on the
% left; on the right teaching prose the student completes with inline \blank{} fills, every
% multi-step computation in an inline `work` block, then a grid of short numbered problems with
% reserved answer space. Problems are numbered 1..20 through the whole component.
%   I DO   : the teacher narrates the prose, fills the lines, and works the FIRST problem of each
%            grid; every `work` block is teacher-worked at the board.
%   WE DO  : the class works the rest of each grid, students holding the pen; the last row,
%            Guided Practice, is worked entirely together.
%   YOU DO : nothing on this page — ../homework/ IS the individual practice, started in the
%            last minutes of the period. No independent set, no reflection box, no objectives
%            box (the targets are on the cover), no notesbox, no practicebox.
% Vocabulary is GIVEN here, not discovered.
%
% LAYOUT, modelled row-for-row on unit01/lesson05 (the house style):
%   * 5 rows, 20 problems, 5 `work` blocks, 14 sub-row seams, EXACTLY TWO TikZ figures.
%   * The Notes column is teaching prose carrying the fills mid-sentence, never a terse
%     definition followed by a big grid — but TERSE prose: 1.5's row 1 runs about six typeset
%     lines and so does this one. Connective and motivational sentences are cut, not kept.
%   * Every computation lives in a `work` block, so no \ans{} ever carries arithmetic. Maya's and
%     Dev's z-scores share ONE block (two columns of one aligned row) rather than two blocks —
%     each block costs \abovedisplayskip + \belowdisplayskip, so merging buys most of an inch.
%   * Data is a compact \rowcolor{sky} tabularx or a centred number list, not a picture, unless
%     the picture itself is the object of study. The old bread HISTOGRAM became the band-count
%     table; the old gram/z number line was dropped (mu and sigma are stated in the prose).
%   * Each row is split into 3–5 sub-rows (a bare `\\`, then a line beginning `& `) so the
%     table can break between pages; the problem grid is ALWAYS the last sub-row of its row.
%   * Grids hold exactly 3 problems (|Y|Y|Y|); the crux row uses 2x2 and Guided Practice 3x2.
%     \pcell heights stay 0.9–1.2cm.
%   * Axis captions clear their tick labels: at scale 0.50 a caption needs ~1.0 units of
%     clearance below the tick-label row, not 0.55. Panel titles are SHORT and carry the units,
%     so the two dotplots need no bottom captions and their titles cannot touch.
%
% THE CRUX is row 4: that the bigger RAW number is the better result — that a raw score can be
% compared across two different distributions without standardizing (problem 11), with its
% companion error, reading a negative z as automatically worse when in a race lower is better
% (problem 13). Problem 14 is the counterweight: what makes a data set NOT produce 68/95/99.7.
%
% THE PROTECTED DATA (the TikZ is pre-drawn; nothing here is student-drawn):
%   Shot put field (20 throws, m): mean 9.0, s 1.5. Maya 12.0 -> z = (12.0-9.0)/1.5 = +2.0.
%   100 m field (20 times, s): mean 13.2, s 0.6. Dev 12.3 -> z = (12.3-13.2)/0.6 = -1.5,
%     AND LOWER IS BETTER IN A SPRINT.
%   Bread: 500 loaves, mean 800 g, spread 20 g. Bands 11/68/170/171/67/12 (+1 outside):
%     within 1 = 341 (68.2%), within 2 = 476 (95.2%), within 3 = 499 (99.8%).
%   Clinic: 400 readings, mean 98.2 F, spread 0.7 F. Bands 9/54/137/136/54/9 (+1 outside):
%     within 1 = 273 (68.25%), within 2 = 381 (95.25%), within 3 = 399 (99.75%).
%   The two records' percentages agreeing is the second discovery (EK VAR-2.A.3).
%   Guided Practice — the bakery as a model: mu = 800 g, sigma = 20 g. 834 g -> z = 1.7;
%     840 g is +2 s.d., so about 2.5% above it (5% outside, halved by symmetry);
%     754 g -> z = -2.3; 790 g -> z = -0.5.
%
% PAGE LOCKSTEP with ../notes/: H is the ANSWER space below the statement and blank widths are
% sized to the answers, so the two files set to the same page count. Every \pcell answer is kept
% well inside its H, or it overflows silently.
% \ans{} IS TEXT MODE: this lesson is the highest-risk one for it. Never put \ans{} inside
% $...$ or \[...\]; every symbol in an answer is wrapped, as \ans{$\mu$}, \ans{$\bar x$}. The
% arithmetic lives in `work` blocks precisely so that no \ans{} has to carry a formula.
% TABLE MECHANICS: inside a cell, `\\` ENDS THE ROW — break lines with \par. (Inside a `work`
% block `\\` is the amsmath aligned's own row break and is safe.)
% NO SKETCH QUESTIONS: every dotplot and curve is pre-drawn and read.
% Target 3 pages at 12pt: the vocabulary box plus a two-page table.
\documentclass[12pt]{article}
\usepackage{apstats-article}
\usepackage{apstats-key}   % pulls in -boxes; do NOT also load -boxes

\begin{document}

\pageheader{Unit 1, Lesson 1.7}{Guided Notes \& Practice --- Answer Key}
% NAMESTRIP: no \namedateperiod here — the name row lives on the cover only.

\vspace{2pt}

% ── VOCABULARY (I do — students fill each term AS IT IS NAMED, never in advance) ──
\boxguard[16]
\begin{vocabbox}
\small
Fill in each term \textbf{as we name it} in the notes below.
\par
\termans{Parameter}{a numerical summary of a population --- $\mu$ for its mean, $\sigma$ for its standard deviation.}
\par
\termans{Statistic}{a numerical summary of a sample --- $\bar x$ and $s_x$ (Lesson 1.5).}
\par
\termans{Normal distribution}{a symmetric, mound-shaped curve fixed entirely by $\mu$ and $\sigma$.}
\par
\termans{Empirical rule}{in an approximately normal distribution, about 68\%, 95\%, and 99.7\% of observations lie within 1, 2, and 3 standard deviations of the mean.}
\par
\termans{Standardized score ($z$-score)}{$z = (x_i - \mu)/\sigma$ --- how many standard deviations a value sits from its own distribution's mean.}
\par
\end{vocabbox}

\small
\begin{guidednotes}
% ── ROW 1: parameter vs statistic ────────────────────────────────────────────
\mainidea[Parameter \&]{Statistic} &
Ask first: \textbf{who got measured?} A \textbf{parameter} summarizes a \ans{population} ---
\emph{everybody}. Its symbols are \ans{$\mu$} for the mean and \ans{$\sigma$} for the standard
deviation.
\\
& A \textbf{statistic} summarizes a \ans{sample} --- \emph{some of them}. Its symbols are
Lesson~1.5's pair: \ans{$\bar x$} and \ans{$s_x$}. Only the \ans{object} changes.
\par\vspace{2pt}
\textbf{The test:} ask \ans{who got measured} --- everyone, or a handful? All 500 loaves a bakery baked:
\ans{parameter}. The 40 an inspector weighed: \ans{statistic}.
\\
& \notesprompt{Label each summary a parameter (P) or a statistic (S), and give its symbol.}
\begin{probgrid}{|Y|Y|Y|}
\pcell{1}{The mean distance of all 20 throwers in Saturday's shot put field}{1.0cm}{P --- $\mu$} &
\pcell{2}{A coach times 6 of the meet's 20 sprinters and averages those}{1.0cm}{S --- $\bar x$} &
\pcell{3}{The standard deviation of all 400 temperatures a clinic recorded}{1.0cm}{P --- $\sigma$} \\ \hline
\end{probgrid}
\\ \hline
% ── ROW 2: the standardized score — two tryouts ──────────────────────────────
\mainidea[Standardized score]{$z$-score} &
One scholarship, two events: Maya threw 12.0 m; Dev ran 12.3 s. Nearly equal numbers, and not on
the same \ans{scale}. A \textbf{standardized score}, or \textbf{$z$-score}, counts how many
\ans{standard deviations} a value sits from the mean of its \emph{own} distribution:
\par\vspace{2pt}
\centerline{$z = \dfrac{x_i - \mu}{\sigma}$}
\par\vspace{2pt}
\stepnum{1}\ Subtract the mean: a distance in the variable's own \ans{units}.\par
\stepnum{2}\ Divide by $\sigma$: the units \ans{cancel}, leaving a plain count.\par
\stepnum{3}\ Keep the \ans{sign}: positive is \ans{above} the mean, negative \ans{below} it.
\\
& {\centering
\begin{tikzpicture}[scale=0.50]
  % ── left: the shot put field, one dot per throw. Short title carries the units, so the
  %    panel needs no bottom caption and the two titles cannot collide. ──
  \draw[->, thick] (-0.2,0) -- (5.7,0);
  \foreach \v in {6,7,8,9,10,11,12} {
    \draw ({(\v-5.5)*0.8},-0.07) -- ({(\v-5.5)*0.8},0.07);}
  \foreach \v in {6,8,10,12} {
    \node[below, font=\tiny] at ({(\v-5.5)*0.8},-0.10) {\v};}
  \foreach \v/\n in {6.0/1, 7.0/2, 7.5/1, 8.0/3, 8.5/2, 9.0/3, 9.5/2, 10.0/2, 10.5/1, 11.0/1,
    11.5/1} {
    \foreach \k in {1,...,\n} { \fill[navy] ({(\v-5.5)*0.8},{0.24*\k}) circle (0.10); }}
  \fill[goldacc] ({(12.0-5.5)*0.8},0.24) circle (0.12);
  \node[font=\tiny\bfseries, text=goldacc] at ({(12.0-5.5)*0.8},0.60) {Maya};
  \node[font=\bfseries\scriptsize] at (2.7,1.35) {Shot put (metres)};
  % ── right: the 100 m field, one dot per runner ──
  \begin{scope}[xshift=6.6cm]
    \draw[->, thick] (-0.2,0) -- (5.5,0);
    \foreach \v in {12.5, 13.0, 13.5, 14.0, 14.5} {
      \draw ({(\v-12.2)*2.2},-0.07) -- ({(\v-12.2)*2.2},0.07);}
    \foreach \v/\lab in {12.5/12.5, 13.5/13.5, 14.5/14.5} {
      \node[below, font=\tiny] at ({(\v-12.2)*2.2},-0.10) {\lab};}
    \foreach \v/\n in {12.5/1, 12.6/1, 12.7/1, 12.8/2, 12.9/1, 13.0/1, 13.1/1, 13.2/2, 13.3/1,
      13.4/1, 13.5/1, 13.6/2, 13.7/1, 13.9/1, 14.1/1, 14.5/1} {
      \foreach \k in {1,...,\n} { \fill[navy] ({(\v-12.2)*2.2},{0.24*\k}) circle (0.10); }}
    \fill[goldacc] ({(12.3-12.2)*2.2},0.24) circle (0.12);
    \node[font=\tiny\bfseries, text=goldacc] at ({(12.3-12.2)*2.2},0.60) {Dev};
    \node[font=\bfseries\scriptsize] at (2.65,1.35) {100 m dash (seconds)};
  \end{scope}
\end{tikzpicture}
\par}
\vspace{2pt}
\footnotesize\setlength{\tabcolsep}{3pt}%
\renewcommand{\arraystretch}{1.1}
\begin{tabularx}{\linewidth}{@{} l Y Y Y @{}}
\rowcolor{sky}
\textbf{Event} & \textbf{Field mean} & \textbf{Field std.\ dev.} & \textbf{Our athlete} \\
\hline
Shot put, 20 throwers (metres) & 9.0 & 1.5 & Maya: 12.0 \\
100 m dash, 20 runners (seconds) & 13.2 & 0.6 & Dev: 12.3 \\
\end{tabularx}\small
\\
& Standardize each against \emph{their own} field:
\begin{work}
z_{\text{Maya}} &= \tfrac{12.0 - 9.0}{1.5} = +2.0 \qquad
z_{\text{Dev}} = \tfrac{12.3 - 13.2}{0.6} = -1.5
\end{work}
The units are gone. Maya sits \ans{2.0} standard deviations \ans{above} her field's mean;
Dev sits 1.5 \ans{below} his. \textbf{In a sprint a lower time is better}, so Dev's negative
$z$ is \ans{good} news.
\\
& \notesprompt{Read the two $z$-scores above. Answer in context, sign included.}
\begin{probgrid}{|Y|Y|Y|}
\pcell{4}{Interpret Maya's $z$-score in one sentence, in context.}{1.2cm}{2.0 s.d.\ above her field's mean throw of 9.0 m} &
\pcell{5}{Interpret Dev's. In a sprint, is below average good news?}{1.2cm}{1.5 s.d.\ below the mean 13.2 s --- good in a sprint} &
\pcell{6}{Jordan clears exactly the high-jump field's mean height. Jordan's $z$?}{1.2cm}{$z = 0$ --- exactly at the mean} \\ \hline
\end{probgrid}
\\ \hline
% ── ROW 3: the normal distribution and the empirical rule ────────────────────
\mainidea[Normal distribution \&]{Empirical rule} &
A \textbf{normal distribution} is a curve \ans{symmetric} about its centre and
\ans{mound}-shaped, fixed entirely by \ans{$\mu$} and \ans{$\sigma$}. Real data is never
\emph{exactly} that, so we say a data set is \ans{approximately} normal.
\\
& {\centering
\begin{tikzpicture}[scale=0.50]
  \draw[navy, very thick] plot[smooth] coordinates
    {(-3.5,0.01) (-3,0.03) (-2.5,0.11) (-2,0.34) (-1.5,0.81) (-1,1.52) (-0.5,2.21) (0,2.50)
     (0.5,2.21) (1,1.52) (1.5,0.81) (2,0.34) (2.5,0.11) (3,0.03) (3.5,0.01)};
  \foreach \x/\h in {-3/0.03, -2/0.34, -1/1.52, 0/2.50, 1/1.52, 2/0.34, 3/0.03} {
    \draw[linegray, thin, dashed] (\x,0) -- (\x,\h);}
  \draw[->, thick] (-4.0,0) -- (4.0,0);
  \foreach \x/\lab in {-3/$-3$, -2/$-2$, -1/$-1$, 0/$\mu$, 1/$+1$, 2/$+2$, 3/$+3$} {
    \draw (\x,-0.08) -- (\x,0.08); \node[below, font=\scriptsize] at (\x,-0.12) {\lab};}
  % Caption clearance: tick labels are anchored below at y = -0.12, so at scale 0.50 the
  % caption needs about 1.0 units beneath them. -1.55 leaves 1.43 — deliberately generous.
  \node[font=\scriptsize] at (0,-1.55) {standard deviations from the mean};
\end{tikzpicture}
\par}
\vspace{2pt}
\textbf{Empirical rule.} If a distribution is approximately normal, about \ans{68\%} of
observations lie within 1 standard deviation of the mean, \ans{95\%} within 2, \ans{99.7\%}
within 3.
\\
& Two unrelated records --- a bakery's \textbf{500} loaf weights (mean 800 g, s.d.\ 20 g) and a
clinic's \textbf{400} body temperatures (mean 98.2$^\circ$F, s.d.\ 0.7$^\circ$F) --- each sorted
into bands one standard deviation wide. One loaf and one reading fell outside.
\par\vspace{2pt}
\footnotesize\setlength{\tabcolsep}{2pt}%
\renewcommand{\arraystretch}{1.1}
\begin{tabularx}{\linewidth}{@{} l Y Y Y Y Y Y @{}}
\rowcolor{sky}
\textbf{Band} & $-3$ to $-2$ & $-2$ to $-1$ & $-1$ to 0 & 0 to $+1$ & $+1$ to $+2$ &
  $+2$ to $+3$ \\
\hline
Loaves (500) & 11 & 68 & 170 & 171 & 67 & 12 \\
Readings (400) & 9 & 54 & 137 & 136 & 54 & 9 \\
\end{tabularx}\small
\par\vspace{2pt}
Add the loaves' bands outward:
\begin{work}
\text{within } 1\sigma &: 170 + 171 = 341 \to \tfrac{341}{500} = 68.2\% \\
\text{within } 2\sigma &: 341 + 68 + 67 = 476 \to \tfrac{476}{500} = 95.2\% \\
\text{within } 3\sigma &: 476 + 11 + 12 = 499 \to \tfrac{499}{500} = 99.8\%
\end{work}
\\
& \textbf{7.} Now the clinic's 400 readings, the same way:
\par\vspace{2pt}
\footnotesize\setlength{\tabcolsep}{2pt}%
\renewcommand{\arraystretch}{1.25}
\begin{tabularx}{\linewidth}{@{} l Y Y Y Y @{}}
\rowcolor{sky}
\textbf{Within \ldots{} of the mean} & \textbf{Loaves} & \textbf{Percent of 500} &
\textbf{Readings} & \textbf{Percent of 400} \\
\hline
1 standard deviation & 341 & 68.2\% & \ans{273} & \ans{68.25\%} \\
2 standard deviations & 476 & 95.2\% & \ans{381} & \ans{95.25\%} \\
3 standard deviations & 499 & 99.8\% & \ans{399} & \ans{99.75\%} \\
\end{tabularx}\small
\\
& \notesprompt{Compare your two columns of percentages.}
\begin{probgrid}{|Y|Y|Y|}
\pcell{8}{Grams and degrees, 500 and 400. Why did the percentages come out the same?}{1.2cm}{Shape fixes them --- not the units, not the count} &
\pcell{9}{Round all six percentages into the empirical rule's three numbers.}{1.2cm}{About 68\%, 95\%, and 99.7\%} &
\pcell{10}{Why do we say a data set is only \emph{approximately} normal?}{1.2cm}{Real data is never exactly symmetric --- it is a model} \\ \hline
\end{probgrid}
\\ \hline
% ── ROW 4: THE CRUX — the raw-number comparison, and the sign read wrongly ───
\mainidea[Whose day was]{Better?} &
\textbf{A coach settles the scholarship in one line: ``Dev ran 12.3 and Maya threw 12.0, so Dev's
number is bigger.''} Which is the argument --- that, or the two $z$-scores?
\par\vspace{2pt}
\fcolorbox{navy}{sky}{\parbox{\dimexpr\linewidth-2\fboxsep-2\fboxrule\relax}{\centering
\textbf{A raw value cannot be compared across two different \ans{distributions}.}\par
Standardize each against its own mean and standard deviation, then read the
\textbf{\ans{direction}} in context.}}
\\
& The sign is not a verdict: it says which \ans{side} of the mean the value fell on. A negative
$z$ means \ans{below} the mean, nothing more. Good or bad depends on the \ans{variable}: below
average on a 100 m \emph{time} means \ans{faster} than the field.
\par\vspace{2pt}
Likewise the empirical rule is a fact about a \ans{shape}: a strongly \ans{skewed} or
two-peaked record will not give 68--95--99.7.
\\
& \notesprompt{Settle it properly, then repair two wrong readings.}
\begin{probgrid}{|Y|Y|}
\pcell{11}{\textbf{The coach's argument.} Name what is wrong with comparing 12.0 and 12.3
directly.}{1.2cm}{Metres against seconds --- two fields, never one scale} &
\pcell{12}{Settle it with the two $z$-scores, and say who had the better day \emph{for their own
event}.}{1.2cm}{Maya $+2.0$, Dev $-1.5$ (fast); Maya sits further out} \\ \hline
\pcell{13}{``Dev's $z$ is negative, so Dev ran worse than average.'' Repair that
sentence.}{1.2cm}{Below the mean \emph{time} --- Dev ran faster, so better} &
\pcell{14}{Describe a data set whose bands would \emph{not} come out 68--95--99.7, and say
why.}{1.2cm}{Strongly skewed or two-peaked: the rule needs normal shape} \\ \hline
\end{probgrid}
\\ \hline
% ── LAST ROW: GUIDED PRACTICE — we do, together ──────────────────────────────
\mainidea[Guided practice]{The Bread Loaf} &
The bakery as a \emph{model}: loaf weights are approximately normal, $\mu = 800$ g and
$\sigma = 20$ g. A loaf comes off the line at 834 g:
\begin{work}
z &= \tfrac{834 - 800}{20} = \tfrac{34}{20} = 1.7
\end{work}
That loaf is 1.7 standard deviations \ans{above} the mean weight of \ans{800} grams.
\\
& \textbf{A tail.} About 95\% lie within 2 standard deviations, so about \ans{5\%} lie
outside; the curve is \ans{symmetric}, so that splits evenly:
\begin{work}
840 \text{ g} &\to z = +2, \quad 5\% \text{ outside} \to \tfrac{5\%}{2} = 2.5\% \text{ above}
\end{work}
\\
& \textbf{A complaint.} A customer returns a 754 g loaf as ``way underweight'':
\begin{work}
z &= \tfrac{754 - 800}{20} = \tfrac{-46}{20} = -2.3
\end{work}
Fewer than \ans{2.5\%} of loaves are that light: the complaint has \ans{merit}. But the
model is only \ans{approximately} normal, and one loaf says nothing about the \ans{rest}.
\\
& \notesprompt{We work these together.}
\begin{probgrid}{|Y|Y|Y|}
\pcell{15}{Interpret $z = 1.7$ in one sentence, in context.}{1.0cm}{1.7 s.d.\ above the mean loaf weight of 800 g} &
\pcell{16}{About what percent of loaves weigh \emph{more} than 840 g?}{1.0cm}{About 2.5\%} &
\pcell{17}{Explain the halving step in one sentence.}{1.0cm}{5\% sits outside on \emph{both} ends; symmetry halves it} \\ \hline
\pcell{18}{A 790 g loaf: find its $z$, and say which side of the mean it is on.}{1.2cm}{$z = -0.5$ --- half a standard deviation below} &
\pcell{19}{Is the 754 g complaint fair? Answer with the evidence.}{1.2cm}{Fair: $z = -2.3$, under 2.5\% of loaves are that light} &
\pcell{20}{Name one limit of the model you just judged with.}{1.2cm}{Normal is only approximate, and one loaf is not evidence} \\ \hline
\end{probgrid}
\\ \hline
\end{guidednotes}

% ── THE NOTES END AT GUIDED PRACTICE ─────────────────────────────────────────
% There is deliberately no "you do" here: ../homework/ is the individual practice,
% started in class in the last 8 minutes while the teacher circulates.

\end{document}
